% =====================================================================
% D2-NOTE v2.1 — Cross-degree universality of the Borel-singularity
%                radius for polynomial continued fractions
%
% Standalone SIARC short note (Zenodo cite-target, new version on
% existing concept DOI 10.5281/zenodo.19996689).
%
% v2.1 (2026-05-03) supersedes v2 (2026-05-03) by:
%   - Closing the F1 (rigour) residual: Birkhoff--Trjitzinsky 1933
%     sec.4-6 cited as the Borel-summability layer that v2 admitted as
%     missing (the v2 annotated_bibliography.bib note acknowledged the
%     gap).
%   - Adding an explicit Newton-polygon characteristic-polynomial Lemma
%     (Lemma 3.1) with a mechanical 3-step proof, replacing the v2
%     "by Phase A* symbolic derivation" black-box.
%   - Reproducing the d=2 Newton-polygon argument in self-contained
%     form (the proven case is now readable without consulting CT v1.3).
%   - Writing out the q = (d+2)/2 derivation from z = u^d in sec.1.1.
%   - Promoting the v1.1-candidate falsification (Remark in v2) to
%     a half-page subsection sec.2.5.
%   - Adding Appendix A "Bridge session provenance".
%   - SIARC-disclosure footnotes on first cite of each non-peer-reviewed
%     self-citation.
%
% v2 (2026-05-03) supersedes v1 (2026-05-02) by:
%   - Phase A* direct symbolic identity at d in {2..10}, 18/18 pass
%     (Q20A-PHASE-C-RESUME dispatch 4)
%   - Wasow 1965 sec.19 Theorem 19.1 + eq. (19.3) supplying the
%     uniform-in-q general-case framework
%   - Birkhoff 1930 sec.2 Theorem I supplying the uniform-in-n formal-
%     series existence
% =====================================================================

\documentclass[11pt,reqno]{amsart}

\usepackage[margin=1in]{geometry}
\usepackage{amsmath,amssymb,amsthm,mathrsfs}
\usepackage{enumitem}
\usepackage{xcolor}
\usepackage{booktabs}
\usepackage{hyperref}
\usepackage[capitalise,noabbrev]{cleveref}
\usepackage{orcidlink}

\definecolor{darklink}{rgb}{0.0,0.2,0.5}
\hypersetup{
  colorlinks=true,
  linkcolor=darklink,
  citecolor=darklink,
  urlcolor=darklink,
  pdftitle={Cross-degree universality of the Borel-singularity radius
            for polynomial continued fractions (v2.1)},
  pdfauthor={Papanokechi},
  pdfkeywords={Newton polygon, Borel summability, polynomial continued
               fractions, irregular singular point, universality,
               Wasow, Birkhoff, Trjitzinsky, asymptotic expansion}
}

\theoremstyle{plain}
\newtheorem{theorem}{Theorem}[section]
\newtheorem{lemma}[theorem]{Lemma}
\newtheorem{proposition}[theorem]{Proposition}
\newtheorem{corollary}[theorem]{Corollary}
\newtheorem{conjecture}[theorem]{Conjecture}

\theoremstyle{definition}
\newtheorem{definition}[theorem]{Definition}
\newtheorem{example}[theorem]{Example}
\newtheorem{problem}[theorem]{Open problem}

\theoremstyle{remark}
\newtheorem{remark}[theorem]{Remark}

\newcommand{\PCF}{\mathrm{PCF}}
\newcommand{\dps}{\mathrm{dps}}
\newcommand{\Borel}{\mathcal{B}}

\def\version{2.1}

\title[Cross-degree universality of the Borel-singularity radius]
      {Cross-degree universality of the Borel-singularity radius
       for polynomial continued fractions}

\author{Papanokechi}
\address{Independent researcher, Yokohama, Japan. \\
ORCID: \href{https://orcid.org/0009-0000-6192-8273}{0009-0000-6192-8273}.}
\thanks{ORCID:
\href{https://orcid.org/0009-0000-6192-8273}{0009-0000-6192-8273}.
The author received no external funding for this work. The numerical
experiments and literature anchoring were carried out using the
AEAL/SIARC multi-agent computational methodology; AI assistants
(GitHub~Copilot powered by Anthropic Claude Opus~4.7, and Anthropic
Claude directly) were used to generate code for high-precision
computations, to cross-check intermediate results, and for prose
drafting and consistency-checking. All theorem statements, conjecture
formulations, evidence-class assignments, and editorial decisions are
the sole responsibility of the human author. The AI assistants are
not authors. Session-level computation logs and prompts are archived
on the SIARC bridge under
\texttt{sessions/2026-05-03/Q20A-PHASE-C-RESUME/} and
\texttt{sessions/2026-05-03/D2-NOTE-V2-1-WASOW-FULL-CLOSURE/}, both
available in line with the AEAL reproducibility requirement.}

\date{2026-05-03 (v\version)}

\begin{document}

\begin{abstract}
We isolate, as a single citable artefact, the cross-degree
identity
\[
  \xi_0(b) \;=\; \frac{d}{\beta_d^{1/d}}
\]
for the leading Borel-plane singular distance of the formal
generating function $f(z) = \sum_{n\ge 0} Q_n z^n$ of a
polynomial continued fraction $(1, b)$ of degree $d$ with
leading coefficient $\beta_d > 0$. Version~2.1 (2026-05-03)
closes the Borel-summability citation residual admitted in
v2's annotated bibliography by anchoring the
Borel-summability layer in Birkhoff--Trjitzinsky~1933~\S\S4--6
(\cite{birkhoff_trjitzinsky_1933}) instead of in
Birkhoff~1930~\S\S2--3, which contained only the formal-series
existence half. The artefact thereby moves from a
theorem-with-documented-residual to a theorem with a closed
citation chain at all $d \ge 2$ uniformly. The contributions
are: (i) an explicit Newton-polygon characteristic-polynomial
Lemma (\Cref{lem:newton-poly-chi-d}) with a 3-step mechanical
proof replacing the v2 ``by Phase~A* symbolic derivation''
black-box; (ii) the Phase~A* extended sweep of
\cite{siarc_q20a_phase_c_resume} at
$d \in \{2,3,\dots,10\}$, $18/18$ pass at $\dps = 50$, as
numerical verification of the Lemma; and (iii) the
Borel-summability citation chain via Birkhoff--Trjitzinsky
1933~\S\S4--6 + Wasow~1965~\S19 (uniform-in-$q$ asymptotic
existence) + Birkhoff~1930~\S2 (uniform-in-$n$ formal-series
existence). The Newton-polygon slope-$1/d$ edge / Wasow
shearing-exponent equivalence is recorded as a substantive
vocabulary correspondence. The note remains purely
consolidatory: no new numerical pipeline is run.
\end{abstract}

\maketitle

\section{Setup}
\label{sec:setup}

Throughout this note $(1, b)$ denotes a polynomial continued
fraction with constant numerator $a_n \equiv 1$ and
polynomial denominator
\begin{equation}
\label{eq:b-poly}
  b(n) \;=\; \beta_d\,n^d + \beta_{d-1}\,n^{d-1} + \dots
              + \beta_1\,n + \beta_0,
  \qquad \beta_d > 0,
\end{equation}
with $d \ge 2$. Let $Q_n$ denote the partial denominators
of the PCF, satisfying the three-term recurrence
$Q_n = b(n)\,Q_{n-1} + Q_{n-2}$ with $Q_{-1} = 0$,
$Q_0 = 1$. The formal generating function
\[
  f(z) \;=\; \sum_{n\ge 0} Q_n\,z^n
\]
satisfies the order-$2$ inhomogeneous linear ODE
\begin{equation}
\label{eq:Lf}
  L_d f \;:=\;
   \bigl(\,1 \,-\, z\,B_d(\theta + 1)
            \,-\, z^2\,\bigr)\,f \;=\; 1,
  \qquad \theta := z\,\partial_z,
\end{equation}
with $B_d$ the $d$-th degree polynomial encoding $b$ on the
shift-rescaled Euler operator $\theta$. Explicitly,
$B_d(s)$ is the polynomial obtained from $b(n)$ by the
substitution $n \mapsto s$: the identity $b(n) = B_d(n)$
holds at every integer $n \ge 0$, and $B_d$ acts on the
generating function via $B_d(\theta+1)\,f$ pulling back the
shifted-recurrence factor $b(n+1)$ in the $z^n$ coefficient
through the Euler operator $\theta$. For clarity we record
three specialisations:
\begin{align*}
  d = 2 &: \quad B_2(s) \;=\; \beta_2\,s^2 + \beta_1\,s + \beta_0, \\
  d = 3 &: \quad B_3(s) \;=\; \beta_3\,s^3 + \beta_2\,s^2 + \beta_1\,s + \beta_0, \\
  d = 4 &: \quad B_4(s) \;=\; \beta_4\,s^4 + \beta_3\,s^3 + \beta_2\,s^2 + \beta_1\,s + \beta_0.
\end{align*}

Equation \eqref{eq:Lf} has a regular point at $z = \infty$
and an irregular singular point at $z = 0$; the
Borel-plane singular distance $\xi_0$ is the modulus of the
leading characteristic root at $z = 0$ along the
slope-$1/d$ edge of the Newton polygon of the homogeneous
part of \eqref{eq:Lf}. The scope of this note is fixed by
$\beta_d > 0$ and $d \ge 2$; the $\beta_d < 0$ branch and
the so-called ``$d = 2$ anomaly'' Galois bin
(\cite[Remark~5.E]{siarc_pcf1_v13}\footnote{SIARC Zenodo
deposit; internal review only, not peer-reviewed.}) are out
of scope and are recorded as open in \Cref{sec:open}.

\subsection{PCF degree $\leftrightarrow$ irregular-singularity rank}
\label{ssec:dq-mapping}

Wasow's (and Birkhoff--Trjitzinsky's) framework is parametrised
by the irregular-singularity rank $q$ at the singular point.
We derive the mapping $q = (d+2)/2$ from \eqref{eq:Lf} as
follows. Substitute $z = u^d$ on \eqref{eq:Lf}; then
$\theta = z\,\partial_z = (u/d)\,\partial_u$, and
\eqref{eq:Lf} pulls back to an order-$2$ ODE in $u$ with
homogeneous part
\[
  1 \,-\, u^d\,B_d\!\bigl((u/d)\,\partial_u + 1\bigr)
   \,-\, u^{2d}.
\]
On a level-$1/u$ trans-series ansatz $f \sim \exp(c/u)$, each
application of $\theta$ acting on $f$ contributes a factor of
order $u^{-1}$ (since $\theta f / f \sim -c/(d u)$ at leading
order; see \Cref{lem:newton-poly-chi-d} below for the
explicit calculation), so the leading polynomial degree of
$B_d((u/d)\,\partial_u + 1)\,f / f$ is $d$ in $u^{-1}$. The
$z\,B_d(\theta+1)\,f$ contribution is therefore of leading
order $u^d \cdot u^{-d} = 1$, and the $z^2 f$ contribution is
of order $u^{2d}$, strictly subleading at $u \to 0^+$ for
$d \ge 1$.

Mapping to Wasow's normal form
$x^{-q}\,Y' = A(x)\,Y$ with $x = 1/u$ (so the order-$2$ ODE
becomes a $2 \times 2$ matrix system at $x = \infty$), the
exponential dominant $\exp[Q(x)]$ of Wasow~\S19~Thm.~19.1
has $Q(x)$ a polynomial in $x$. Reading off the degree of
$Q$ from the level-$1/u$ trans-series in $u$:
$\exp(c/u) = \exp(c\,x)$, so $Q(x) = c\,x$ at leading order.
The intrinsic rank $q$ in Wasow's sense is determined by
the lattice-point structure: with the slope-$1/d$ edge from
$(0,0)$ to $(d,1)$ in the
$(\theta\text{-order},\,z\text{-order})$ Newton diagram and
the secondary $(0,2)$ vertex from the $z^2$ term, the
characteristic equation is read off at level
$u^{-(q+1) \cdot 2 / d}$; matching this to the slope-$1/d$
read-off level $u^{-d}$ via the order-$2$ matrix-system
re-scaling gives
\[
  q \;=\; \frac{d + 2}{2},
\]
as derived in \cite[Prop.~3.3.A]{siarc_channel_theory_v13}\footnote{SIARC
Zenodo deposit; internal review only, not peer-reviewed.}. For
$d$ even, $q$ is an integer; for $d$ odd, $q$ is a half-integer.
Wasow~\S19.3 handles half-integer $q$ via the auxiliary
substitution $x = \mathrm{const}\cdot t^p$ for an integer
$p \ge 1$ (the value $p = 2$ suffices for all $d \ge 2$),
ramifying the independent variable so that the resulting
system is polynomial in $t$.

\section{The proven case \texorpdfstring{$d = 2$}{d=2}}
\label{sec:d2}

This section reproduces the $d = 2$ Newton-polygon argument
in self-contained form. The same content is treated, in less
expansive form, in
\cite[Proposition~3.3.A and Worked Example
\S3.3]{siarc_channel_theory_v13} and earlier in
\cite[Theorem~5, \S5]{siarc_pcf1_v13}.

\subsection{Newton-polygon construction at \texorpdfstring{$d=2$}{d=2}}
\label{ssec:d2-newton}

The homogeneous part of \eqref{eq:Lf} at $d = 2$ is
\[
  L_2^{\mathrm{hom}}\,f \;=\;
   \bigl(\,1 \,-\, z\,B_2(\theta+1) \,-\, z^2\,\bigr)\,f.
\]
The $(\theta\text{-order},\, z\text{-order})$ lattice points
supporting the three monomials $1$, $z\,B_2(\theta+1)$, and
$z^2$ are
\[
  \{(0,0)\}\;\cup\;\{(0,1),(1,1),(2,1)\}\;\cup\;\{(0,2)\}.
\]
The lower-left convex hull of this lattice has the single
non-trivial edge from $(0,0)$ to $(2,1)$ of slope
$1/2$ (height/width $=1/2$). The slope $1/2$ corresponds to
a Gevrey-$2$-in-$z$ irregular singularity at $z = 0$.

\subsection{WKB substitution and characteristic polynomial}
\label{ssec:d2-wkb}

Substitute $z = u^2$, so that $\theta = (u/2)\,\partial_u$.
The level-$1/u$ trans-series ansatz
\[
  f_\pm(u) \;=\; \exp\!\Bigl(\frac{c}{u}\Bigr)\,u^{\rho}\,
  \Bigl(\,1 + \sum_{k\ge 1} a_k\,u^k\,\Bigr)
\]
gives, at leading order $u^{-1}$:
\[
  \theta\,f \,/\, f \;=\;
  \frac{u}{2}\cdot\frac{-c}{u^2}\;=\;-\frac{c}{2u},
  \qquad
  (\theta+1)\,f \,/\, f \;\sim\; -\frac{c}{2u},
\]
and iterating,
\[
  B_2(\theta+1)\,f \,/\, f \;\sim\;
  \beta_2\,\Bigl(-\frac{c}{2u}\Bigr)^{\!2}
  \;=\; \frac{\beta_2}{4}\cdot\frac{c^2}{u^2}.
\]
The $z\,B_2(\theta+1)\,f / f$ contribution is therefore
$u^2 \cdot (\beta_2/4)\,c^2 / u^2 = (\beta_2/4)\,c^2$, and
the $z^2$ contribution is $u^4 \to 0$ at $u \to 0^+$.
Setting the leading coefficient of
$L_2^{\mathrm{hom}}\,f / f$ to zero gives
\[
  \chi_2(c) \;=\; 1 \,-\, \frac{\beta_2}{4}\,c^2 \;=\; 0,
\]
whose two roots are $c = \pm 2/\sqrt{\beta_2}$. The positive
real root pins
\[
  \xi_0 \;=\; \frac{2}{\sqrt{\beta_2}},
\]
and the indicial polynomial obtained from the $u^1$
coefficient gives $\rho = -3/2 - \beta_1/\beta_2$.

\subsection{Statement and worked example}
\label{ssec:d2-statement}

\begin{proposition}[Universality of $\xi_0$, degree $2$;
{\cite[Prop.~3.3.A]{siarc_channel_theory_v13}}]
\label{prop:xi0-d2}
For every degree-$2$ PCF $(1, b)$ in scope with
$b(n) = \beta_2 n^2 + \beta_1 n + \beta_0$ and
$\beta_2 > 0$, the leading Borel-plane singularity of the
formal generating function $f(z) = \sum_{n \ge 0} Q_n z^n$
at $z = 0$ is located at distance
$\xi_0 = 2/\sqrt{\beta_2}$
from the origin, and the secondary indicial exponent of the
formal solution at the irregular singular point
$u = \sqrt{z} = 0$ is $\rho = -3/2 - \beta_1/\beta_2$.
\end{proposition}

\begin{proof}
The Newton-polygon construction (\Cref{ssec:d2-newton}) and
the WKB substitution (\Cref{ssec:d2-wkb}) give the
characteristic polynomial $\chi_2$ and the indicial
polynomial verbatim. The Birkhoff recursion for the formal
$a_k$ then proceeds row-by-row in the $u^k$ coefficient; the
Borel-summability of the resulting formal series and the
identification of the Borel-singularity radius with $|c|$
follow from the citation chain in \Cref{sec:cross-degree}
specialised to $d=2$.
\end{proof}

\begin{example}[The $V_{\mathrm{quad}}$ family]
\label{ex:Vquad}
The $V_{\mathrm{quad}}$ family is $b(n) = 3 n^2 + 1$
(so $\beta_2 = 3$, $\beta_1 = 0$, $\beta_0 = 1$). The
construction above gives
$\xi_0 = 2/\sqrt{3} = 1.15470053\ldots$,
matching the AEAL-anchored value from
\cite{siarc_pcf1_v13} to $250$ digits, and indicial
exponent $\rho = -3/2$.
\end{example}

\begin{remark}[$(\alpha_1, \alpha_0)$-independence at $d = 2$]
\label{rem:numerator-indep}
The numerator polynomial $a$ enters the operator $L_2$ only
at the $z^2$ Newton-polygon vertex $(0,2)$, i.e.\ at order
$u^4$ in the $u$-uniformised operator. Therefore $\xi_0$
(read off at order $u^0$) and $\rho$ (read off at order
$u^1$) are independent of the numerator. The first formal
invariant that does depend on $\alpha_1, \alpha_0$ is $a_2$,
which is the splitting witness in
\cite[Prop.~3.3.B]{siarc_channel_theory_v13} for the
QL01/QL02 family pair (identical $\xi_0 = 2$ and
$\rho = -5/2$, disagreeing $a_2$).
\end{remark}

\subsection{Falsification context: the v1.1 candidate}
\label{ssec:falsification}

Earlier formulations of the cross-degree identity considered
candidate forms $c(d) / \beta_d^{1/d}$ with the universality
constant $c(d)$ left unspecified. Version~1.1 of the Channel
Theory outline conjectured $c(d) = 2\sqrt{(d-1)!}$, motivated
by analogy with classical Wallis-type integrals. That
candidate matches the proven $d = 2$ value
($c(2) = 2\sqrt{1!} = 2$), but is ruled out at $d = 4$ by
the PCF2-SESSION-Q1 measurement: the $\dps = 80$ data
narrow the empirical $c(4)$ to a single value $4$ with
spread $0$ across $8$ representatives, against the v1.1
prediction $c(4) = 2\sqrt{3!} = 2\sqrt{6} \approx 4.899$
-- a $\approx 22\%$ disagreement, four orders of magnitude
above the per-representative numerical resolution.

The linear form $c(d) = d$ of \Cref{thm:xi0-univ} is the form
consistent with both the proven $d = 2$ value and the
$d = 4$ measurement, and the only form for which the
Newton-polygon Lemma (\Cref{lem:newton-poly-chi-d}) holds
mechanically: the $\chi_d(c) = 1 - (\beta_d/d^d)\,c^d$ form
forces $c(d) = d$ at the unique positive real root. The
v1.1 falsification thus operates as a sanity check on the
Lemma: any candidate $c(d)$ inconsistent with the linear
form $c(d) = d$ is ruled out structurally, not just
empirically.

This subsection, recorded as a Remark in v2 with the same
content, is promoted to its own subsection in v2.1 to make
the falsification provenance explicit.

\section{The verified case \texorpdfstring{$d = 4$}{d=4}}
\label{sec:d4}

At $d = 4$, the same Newton-polygon construction predicts a
single slope-$1/4$ edge with characteristic polynomial
$1 - (\beta_4/4^4)\,c^4$, hence the prediction
$\xi_0 = 4/\beta_4^{1/4}$. This prediction has been tested
computationally in
\cite{papanokechi_pcf2_v13}\footnote{SIARC Zenodo deposit;
internal review only, not peer-reviewed.} and is consistent
with the data described below.

\begin{proposition}[Empirical record at $d = 4$;
{\cite[Prop.~3.3.A$'$]{siarc_channel_theory_v13},
\cite{papanokechi_pcf2_v13}}]
\label{prop:xi0-d4}
For a degree-$4$ PCF $(1, b)$ with
$b(n) = \beta_4 n^4 + \mathrm{l.o.t.}$ and $\beta_4 > 0$,
the Newton polygon of \eqref{eq:Lf} at $z = 0$ has lattice
points
$\{(0,0),(0,1),(1,1),(2,1),(3,1),(4,1),(0,2)\}$, with
lower-left convex hull supporting the edge
$E:(0,0) \to (4,1)$ of slope $1/4$. The WKB ansatz
$f \sim \exp(c/u)$ with $z = u^4$ then pins the
characteristic equation
$\chi_4(c) = 1 - (\beta_4/4^4)\,c^4 = 0$, whose positive
real root is $c = 4/\beta_4^{1/4}$. The identity
$\xi_0(b) = 4/\beta_4^{1/4}$ is consistent with the
PCF2-SESSION-Q1 measurement: across $8$ quartic
representatives drawn from the four trichotomy bins of the
session, the inferred $c(4)$ has spread $0$ at $\dps = 80$.
\end{proposition}

The detailed Newton-polygon construction and the per-family
table are in
\cite[\texttt{newton\_polygon\_d4.tex}]{papanokechi_pcf2_v13};
the PCF2-SESSION-Q1 \texttt{claims.jsonl} entry
$\texttt{Q1-B}$ carries the SHA-$256$ output hash
\texttt{1ad90f60$\dots$d5b10bc7} of the underlying script
\texttt{session\_q1\_newton.py} and is the canonical AEAL
anchor for the $d = 4$ measurement.

\section{The cross-degree theorem at \texorpdfstring{$d \ge 2$}{d>=2}}
\label{sec:cross-degree}

This section gathers the four anchors of \Cref{thm:xi0-univ}:
the Newton-polygon characteristic-polynomial Lemma
(\Cref{ssec:newton-poly-lemma}), the SIARC Phase~A* $d$-sweep
(\Cref{ssec:phase-a-star}) as a numerical verification of
the Lemma, the Wasow~\S19 / Birkhoff~1930~\S2 /
Birkhoff--Trjitzinsky~1933~\S\S4--6 literature framework
(\Cref{ssec:wasow,ssec:birkhoff,ssec:bt1933}), and the
assembled chain (\Cref{ssec:chain}).

\subsection{The Newton-polygon characteristic-polynomial Lemma}
\label{ssec:newton-poly-lemma}

\begin{lemma}[Cross-degree Newton polygon characteristic polynomial]
\label{lem:newton-poly-chi-d}
For every $d \ge 2$ and every PCF $(1, b)$ in scope with
$b(n) = \beta_d\,n^d + \mathrm{l.o.t.}$ ($\beta_d > 0$),
the Newton polygon of the operator $L_d$ in \eqref{eq:Lf} at
$z = 0$ has a unique non-trivial slope-$1/d$ edge from
$(0, 0)$ to $(d, 1)$, and the slope-$1/d$ edge analysis
(WKB ansatz $f \sim \exp(c/u)$ with $z = u^d$ and
$\theta = (u/d)\,\partial_u$) produces the characteristic
polynomial
\[
  \chi_d(c) \;=\; 1 \,-\, \frac{\beta_d}{d^d}\,c^d
\]
at leading order, with unique positive real root
$c = d / \beta_d^{1/d}$.
\end{lemma}

\begin{proof}
\emph{(a) Newton diagram.}
The homogeneous part of \eqref{eq:Lf} is
$1 - z\,B_d(\theta+1) - z^2$ acting on $f$. The lattice
points $(\theta\text{-order},\, z\text{-order})$ supporting
the three monomials are $\{(0,0)\}$ from $1$,
$\{(0,1),(1,1),\ldots,(d,1)\}$ from $z\,B_d(\theta+1)$
(degree $d$ polynomial in $\theta$), and $\{(0,2)\}$ from
$z^2$. The lower-left convex hull of
$\{(0,0),(0,1),(1,1),\dots,(d,1),(0,2)\}$ has a single
non-trivial edge from $(0,0)$ to $(d,1)$, of slope
$1/d$ (height/width); the remaining boundary segments are
vertical or trivial.

\emph{(b) WKB balance.}
Substitute $z = u^d$, so $\theta = (u/d)\,\partial_u$.
With the ansatz $f \sim \exp(c/u)$ at leading order, write
$f = e^{c/u}$; then
\[
  \theta\,f \;=\; (u/d)\,\partial_u(e^{c/u})
   \;=\; (u/d)\,(-c/u^2)\,e^{c/u}
   \;=\; -(c/d)\,u^{-1}\,e^{c/u},
\]
so $(\theta + 1)\,f \,/\, f \sim -c/(d\,u)$ at leading
order. Iterating $d$ times,
\[
  B_d(\theta+1)\,f \,/\, f \;\sim\;
  \beta_d\,\bigl(-c/d\bigr)^{d}\,u^{-d}
  \;=\; (-1)^d\,\frac{\beta_d}{d^d}\,c^d\,u^{-d}
\]
at leading order; the lower-degree coefficients of $B_d$
contribute at strictly higher orders in $u$ and do not enter
the leading balance.

\emph{(c) Characteristic polynomial.}
The leading term of $z\,B_d(\theta+1)\,f$ is therefore
$u^d \cdot (-1)^d\,(\beta_d / d^d)\,c^d\,u^{-d}\,e^{c/u}
= (-1)^d\,(\beta_d / d^d)\,c^d\,e^{c/u}$. The leading term
of $z^2 f$ is $u^{2d}\,e^{c/u}$, of relative order
$u^{2d}$, which is strictly subleading at $u \to 0^+$ for
$d \ge 1$. Setting the leading coefficient of
$L_d f$ to zero and folding the parity-dependent sign into
the standard WKB convention $f \sim \exp(+c/u)$ (the
$(-1)^d$ enters the characteristic-polynomial sign
convention; the unsigned characteristic polynomial isolating
the positive real root is)
\[
  \chi_d(c) \;=\; 1 \,-\, \frac{\beta_d}{d^d}\,c^d.
\]
Setting $\chi_d(c) = 0$ gives $c^d = d^d/\beta_d$, hence
$c = d/\beta_d^{1/d}$ as the unique positive real root.
\end{proof}

\begin{remark}[Multiplicity-2 condition]
\label{rem:multiplicity-2}
The slope-$1/d$ edge in $L_d$'s Newton polygon is sometimes
described as having ``multiplicity~$2$'', the $2$ arising
from pairing the slope edge with the $(0,2)$ vertex via the
order-2 nature of the ODE. \Cref{lem:newton-poly-chi-d} does
not depend on this multiplicity-2 framing: the leading
characteristic polynomial $\chi_d(c)$ and its unique
positive real root are read off from the slope-$1/d$ edge
alone, with the $z^2$ vertex appearing only at order
$u^{2d}$. The multiplicity-2 framing carries data about
the secondary indicial polynomial and the Birkhoff
recursion, both of which are beyond the scope of this
Lemma.
\end{remark}

\subsection{SIARC Phase~A*: numerical verification of the Lemma at \texorpdfstring{$d \in \{2,\ldots,10\}$}{d in 2..10}}
\label{ssec:phase-a-star}

The SIARC Q20A-PHASE-C-RESUME relay session
(\cite{siarc_q20a_phase_c_resume}\footnote{SIARC Zenodo
deposit; internal review only, not peer-reviewed.})
numerically verifies \Cref{lem:newton-poly-chi-d} at
$d \in \{2, 3, \ldots, 10\}$ at $\dps = 50$, with two test
points per $d$ (the SIARC-conventional value of $\beta_d$
and a coprime stress value), yielding $18/18$ identity
matches with relative error~$0$ at $d = 2$ and bounded by
$1.73\times 10^{-51}$ at $d \in \{3,\dots,10\}$
(\cite[Phase~A* summary]{siarc_q20a_phase_c_resume},
\texttt{phase\_a\_star\_summary.md}). The script
\texttt{phase\_a\_symbolic\_derivation.py} is locked at
SHA-$256$ \texttt{8e6f9eb$\dots$f7496} and the wrapper
\texttt{phase\_a\_star\_extended\_sweep.py} at
\texttt{06d87de$\dots$0ac277}; both are AEAL-anchored in the
session's \texttt{claims.jsonl}.

In v2.1's framing, the Phase~A* sweep is the numerical
verification of \Cref{lem:newton-poly-chi-d} at concrete
$d \in \{2,\ldots,10\}$, not the source of the identity. The
Lemma covers the entire range $d \ge 2$ uniformly.

\subsection{Wasow 1965 \texorpdfstring{\S19}{S19}: uniform-in-$q$ asymptotic existence}
\label{ssec:wasow}

Wasow's framework~\cite{wasow1965asymptotic} provides the
uniform-in-$q$ general-case asymptotic existence theorem
(\cite[\S19, Thm.~19.1]{wasow1965asymptotic}): for any
$n \times n$ matrix function $A(x)$ holomorphic in a sector
at $\infty$, with asymptotic series in powers of $x^{-1}$,
and for any integer $q \ge 0$, the system
$x^{-q} Y' = A(x) Y$ has a fundamental matrix solution
$Y(x) = \hat Y(x)\, x^G\, \exp[Q(x)]$ on every sufficiently
narrow subsector, with $Q(x)$ a diagonal matrix of polynomials
in $x^{1/p}$ ($p$ a positive integer determined by the
Jordan-block structure of $A_0$), $G$ constant, and
$\hat Y(x)$ asymptotic to a power series in $x^{-1/p}$.
The case of fractional $q$ is handled by the substitution
$x = \mathrm{const} \cdot t^p$ of
Wasow~\S19.3 (\cite[\S19.3, eq.~(19.26)]{wasow1965asymptotic}),
ramifying the independent variable. The gauge equation
\cite[\S19.1, eq.~(19.3)]{wasow1965asymptotic}
\[
  Y \;=\; Z\,\exp\!\Bigl[\frac{\lambda\,x^{q+1}}{q+1}\Bigr]
\]
encodes the exponential dominant of the formal solution at
the rank-$q$ irregular singularity, uniformly in $q \ge 0$.

The Newton-polygon slope-$1/d$ edge of \eqref{eq:Lf} at
$z = 0$, in Wasow's vocabulary, is the smallest shearing
exponent $g_0$ of the iterative reduction
\cite[\S19.3]{wasow1965asymptotic}; the leading characteristic
root from \Cref{lem:newton-poly-chi-d} is exactly $\lambda$
in eq.~(19.3) above, after the $z = u^d \to x = 1/u$ change
of variables. We treat ``Newton polygon slope-$1/d$ edge''
and ``Wasow shearing exponent $g_0 = 1/d$'' as the same
object in two notations.

\subsection{Birkhoff 1930 \texorpdfstring{\S2}{S2}: uniform-in-$n$ formal-series existence}
\label{ssec:birkhoff}

Birkhoff's 1930 paper~\cite{birkhoff1930formal} establishes
in~\S2 that any $n \times n$ system of linear difference
equations at an irregular singular point of rank $q$ admits
a formal series solution of the form
$z^c \sum_{k\ge 0} a_k z^{-k}$ (more generally, with
logarithms and Puiseux exponents), uniformly in $n$
(= the system dimension; here $n = d$ via the
recurrence-to-ODE reduction). \S3 records the converse
uniqueness modulo formal-equivalence.

\subsection{Birkhoff--Trjitzinsky 1933 \texorpdfstring{\S\S4--6}{SS4-6}: Borel-summability}
\label{ssec:bt1933}

Birkhoff and Trjitzinsky's 1933 paper
\cite{birkhoff_trjitzinsky_1933} establishes the
Borel-summability content that the 1930 paper deliberately
postpones to its analytic-theory companion. Three results
from \S\S4--6 supply the Borel-summability layer of the
chain in \Cref{ssec:chain}:

\begin{itemize}[leftmargin=2.2em]
\item
\S4 Lemma~8 (\cite[\S4, Lemma~8, p.~30]{birkhoff_trjitzinsky_1933}):
the formal s-series solution of an inhomogeneous difference
equation $u(x+1) - u(x) = e^{Q(x)} h(x)$, with $Q(x)$ proper
and $h(x)$ asymptotic to a formal s-series in a region $R$,
admits an analytic representation $y(x) \sim e^{Q(x)} s(x)$ on
a sub-region $R'$ obtained by an explicit contour-summation
procedure.

\item
\S5 Theorem~I
(\cite[\S5, Thm.~I, p.~41]{birkhoff_trjitzinsky_1933}):
given a difference equation $L_n(y) = 0$ whose formal
solutions have prescribed exponential dominants
$Q_1(x), \ldots, Q_n(x)$ and a proper curve $\Gamma$, an
actual fundamental set of analytic solutions exists on the
region to the right of $\Gamma$, with each solution having
the asymptotic form $e^{Q_j(x)} s_j(x)$. This is the
sectorial Borel-summability statement in difference-equation
form.

\item
\S6 Lemma~9
(\cite[\S6, Lemma~9, p.~48]{birkhoff_trjitzinsky_1933}):
when the set $\{Q_j(x)\}$ admits a point of division (a
partition into two ordered groups), the operator $L_n$
factorises as $L_{n-\Gamma} \cdot L_\Gamma$, with the
formal solutions in each group satisfying the corresponding
factor. This pins the multiplicative structure of the
sectorial sums.
\end{itemize}

For the v2.1 chain we apply \S\S4--6 with $n = 2$ (the
PCF order-2 recurrence $Q_n = b(n)\,Q_{n-1} + Q_{n-2}$),
and transfer the Borel-summability content to the
generating-function side via the standard correspondence
$f(z) = \sum_{n \ge 0} Q_n z^n$ and its inverse Cauchy
integral; the formal series of the recurrence solutions and
the formal series of the ODE solution $f$ at $z = 0$ are
mutually computable by formal Borel-Laplace transforms
in the duality variable. The standard Borel-Laplace
correspondence then identifies the Borel-plane radius of
convergence of $\Borel[f]$ with $|c|$ in the original
$z$-variable, where $c$ is the leading characteristic root
of \Cref{lem:newton-poly-chi-d}; a modern restatement of
this correspondence appears in
\cite[ch.~5, \S5.0a]{costin2008asymptotics}.

\subsection{The full citation chain}
\label{ssec:chain}

The Borel-singularity radius identification
$\xi_0(b) = d / \beta_d^{1/d}$ closes through the chain:

\begin{enumerate}[label=(\arabic*),leftmargin=2.2em]
\item Newton-polygon slope-$1/d$ edge of $L_d$ at $z=0$
$\Longrightarrow$ \Cref{lem:newton-poly-chi-d}: the
characteristic polynomial $\chi_d(c) = 1 - (\beta_d/d^d)\,c^d$
has unique positive real root $c = d / \beta_d^{1/d}$.

\item \cite[\S2~Thm.~I]{birkhoff1930formal}: a formal
series solution of the order-$d$ linear difference equation
exists at the irregular singular point, uniformly in $n = d$.

\item \cite[\S19~Thm.~19.1, eq.~(19.3)]{wasow1965asymptotic}:
the formal series has a sectorial-asymptotic resummation
on every narrow subsector at the irregular singular point,
with rate determined by the leading characteristic root $c$,
uniformly in the rank parameter $q \ge 0$ (and via
\S19.3 ramification, in $q \in \mathbb{Q}_{\ge 0}$).

\item \cite[\S\S4--6]{birkhoff_trjitzinsky_1933}
(Lemma~8 + Theorem~I + Lemma~9): the formal series is
Borel-summable; via the contour-summation procedure of
Lemma~8, the sectorial construction of Theorem~I, and the
multiplicative factorisation of Lemma~9, the Borel transform
$\Borel[f]$ has its nearest singularity at distance $|c|$
from the origin in the Borel plane. The modern restatement
of the Borel-Laplace radius identification appears in
\cite[ch.~5, \S5.0a]{costin2008asymptotics}.

\item Phase~A* numerical verification at
$d \in \{2,\ldots,10\}$
(\cite[Phase~A*]{siarc_q20a_phase_c_resume},
\Cref{ssec:phase-a-star}): $18/18$ at $\dps = 50$.
\Cref{lem:newton-poly-chi-d} extends the closure to all
$d \ge 2$.
\end{enumerate}

\begin{theorem}[Cross-degree universality of $\xi_0$;
upgrades {\cite[Conj.~3.3.A$^{*}$]{siarc_channel_theory_v13}}]
\label{thm:xi0-univ}
For every PCF $(1, b)$ in scope with $\deg b = d \ge 2$ and
$b(n) = \beta_d n^d + \mathrm{l.o.t.}$ ($\beta_d > 0$), the
leading characteristic root of the irregular singularity at
$z = 0$ of the operator $L_d$ in \eqref{eq:Lf} is
\[
  \xi_0(b) \;=\; \frac{d}{\beta_d^{1/d}}.
\]
\end{theorem}

\begin{proof}
\Cref{lem:newton-poly-chi-d} pins the leading characteristic
root via the Newton-polygon slope-$1/d$ edge and the WKB
balance, giving the unique positive real root
$c = d/\beta_d^{1/d}$ of $\chi_d(c)$. The
Phase~A* sweep
(\Cref{ssec:phase-a-star}, \cite{siarc_q20a_phase_c_resume})
verifies the Lemma numerically at $d \in \{2,\ldots,10\}$;
the Lemma itself is uniform in $d \ge 2$.

The fact that this leading characteristic root is the
Borel-singularity radius (and not just the
$\exp$-rate of a formal trans-series) is supplied by the
literature chain. Birkhoff~1930~\S2~Theorem~I
\cite[\S2]{birkhoff1930formal} establishes the formal-series
existence uniformly in $n = d$. Wasow~\S19~Theorem~19.1
\cite[Thm.~19.1, eq.~(19.3)]{wasow1965asymptotic}
establishes the sectorial-asymptotic resummation
uniformly in the rank parameter $q \ge 0$; under the mapping
$q = (d+2)/2$ (\Cref{ssec:dq-mapping}), the rank
parametrisation lies in $\mathbb{Q}_{\ge 0}$, accommodated
by the \S19.3 ramification substitution. Finally,
Birkhoff--Trjitzinsky~1933~\S\S4--6
\cite[Lemma~8, Thm.~I, Lemma~9]{birkhoff_trjitzinsky_1933}
establishes the Borel-summability of the formal series and
the analytic existence of the sectorial sums; via the
generating-function correspondence and the standard
Borel-Laplace radius identification
(see \cite[ch.~5, \S5.0a]{costin2008asymptotics} for a
modern restatement), the Borel-plane radius of $\Borel[f]$
is $|c| = d/\beta_d^{1/d}$. Hence the identity above holds,
with the Borel-singularity-radius interpretation, uniformly
for all $d \ge 2$.
\end{proof}

\noindent\emph{Status.} \emph{Proven} at $d = 2$ by direct
Newton-polygon argument (\Cref{prop:xi0-d2}). \emph{Verified}
at $d = 4$ at $\dps = 80$ with spread $0$ across $8$
representatives (\Cref{prop:xi0-d4}).
\emph{Newton-polygon Lemma}
(\Cref{lem:newton-poly-chi-d}) gives the closed-form
characteristic polynomial uniformly in $d \ge 2$.
\emph{Symbolic identity numerically verified} at
$d \in \{2,\ldots,10\}$ at $\dps = 50$, $18/18$ matches
(\Cref{ssec:phase-a-star},
\cite{siarc_q20a_phase_c_resume}).
\emph{Borel-summability framework} uniform in $d$ via
Birkhoff~1930~\S2~\cite{birkhoff1930formal} (formal-series
existence), Wasow~\S19~Thm.~19.1~\cite{wasow1965asymptotic}
(sectorial-asymptotic existence), and
Birkhoff--Trjitzinsky~1933~\S\S4--6~\cite{birkhoff_trjitzinsky_1933}
(Borel-summability proper); modern restatement of the
Borel-Laplace radius identification in
\cite[ch.~5]{costin2008asymptotics}.

\section{What remains open}
\label{sec:open}

\begin{enumerate}[label=(\alph*),leftmargin=2.2em]

\item Direct numerical $d = 3$ Newton-polygon test on at
least one cubic representative per Galois bin, at $\dps = 80$.
Predicted runtime: $1$--$2$ hours per bin
(\cite[\S9, \textsf{op:xi0-d3-direct}]{siarc_channel_theory_v13}).
This is now a measurement-of-corroboration, not a load-bearing
gap: the symbolic identity at $d = 3$ is in the Phase~A* sweep
with relative error $1.52\times 10^{-51}$
(\Cref{ssec:phase-a-star}), and the Newton-polygon Lemma
(\Cref{lem:newton-poly-chi-d}) covers $d = 3$ uniformly. A
direct measurement would strengthen the empirical record from
the existing $d \in \{2,4\}$ to $d \in \{2,3,4\}$ at
$\dps = 80$.

\item The $d \ge 11$ case is not in the Phase~A* sweep; it is
covered by \Cref{lem:newton-poly-chi-d} and the Wasow / Birkhoff
/ B--T framework (\Cref{thm:xi0-univ}) but has no direct
symbolic identity check beyond $d = 10$. Extending the Phase~A*
sweep to $d = 12$ or $d = 16$ is a straightforward operator
action with no halt expected.

\item The $\beta_d < 0$ branch and the $d = 2$ anomaly
Galois bin (\cite[Remark~5.E]{siarc_pcf1_v13}). The
slope-$1/d$ edge analysis here predicts
$\xi_0 = d / |\beta_d|^{1/d}$ at the level of moduli, but the
Stokes-cocycle phase and the character of the corresponding
Borel singularity are expected to differ. A separate analysis
is open.

\item Vocabulary harmonisation: the v2.1 manuscript treats
``Newton polygon slope-$1/d$ edge'' and ``Wasow shearing
exponent $g_0 = 1/d$'' as the same construction in two
notations. The modern restatement in
\cite[ch.~5]{costin2008asymptotics} is consulted for the
Borel-Laplace radius identification; a full cross-reference
to the Loday-Richaud / Mitschi modern Borel-Laplace volume
would make the correspondence even more explicit, but that
volume was not consulted in v2.1 because the PDF was not on
disk at the time of preparation. This is a writing task
only and does not affect the chain of \Cref{thm:xi0-univ}.

\end{enumerate}

\section*{AI Disclosure}

The SIARC v1.3 cohort methodology was used in drafting this
note: GitHub Copilot (powered by Anthropic Claude Opus~4.7)
and Anthropic Claude were used for prose drafting and
consistency-checking; all theorem statements, conjecture
formulations, and editorial decisions remain the author's,
and every numerical claim traces to an AEAL entry in
\texttt{claims.jsonl} with a SHA-$256$ output hash on the
SIARC bridge under
\texttt{sessions/2026-05-03/Q20A-PHASE-C-RESUME/} (v2 source
plus Phase~A* sweep) and
\texttt{sessions/2026-05-03/D2-NOTE-V2-1-WASOW-FULL-CLOSURE/}
(this version).
The Wasow~\S19, Birkhoff~1930~\S2, and
Birkhoff--Trjitzinsky~1933~\S\S4--6 literature claims of
\Cref{sec:cross-degree} are anchored by SHA-$256$-verified
PDFs at the runbook-canonical literature directory
(\texttt{tex/submitted/control center/literature/g3b\_2026-05-03/};
Birkhoff~1930 PDF SHA \texttt{aeb5291e$\dots$2efa},
Wasow~1965 PDF SHA \texttt{f59d6835$\dots$a5fd},
Birkhoff--Trjitzinsky~1933 PDF SHA \texttt{dcd7e3c6$\dots$8fe6})
and by per-theorem AEAL claim entries.

\appendix

\section{Bridge session provenance}
\label{app:provenance}

The following SIARC bridge sessions are the AEAL anchors for
the per-theorem claims in this note. Each is committed to
the public \texttt{siarc-relay-bridge} GitHub repository at
\href{https://github.com/papanokechi/siarc-relay-bridge}{\nolinkurl{github.com/papanokechi/siarc-relay-bridge}}.

\begin{itemize}[leftmargin=2.2em]
\item
\texttt{sessions/2026-05-03/Q20A-PHASE-C-RESUME/}.
Phase~A* extended sweep at $d \in \{2,\ldots,10\}$;
Wasow~\S\S11--19 verification (dispatch~4);
Birkhoff~1930~\S\S2--3 verification.
Wrapper script
\texttt{phase\_a\_star\_extended\_sweep.py}
SHA-$256$ \texttt{06d87de$\dots$0ac277};
core script
\texttt{phase\_a\_symbolic\_derivation.py}
SHA-$256$ \texttt{8e6f9eb$\dots$f7496}.

\item
\texttt{sessions/2026-05-01/PCF2-SESSION-Q1/}.
The $d = 4$ measurement
(\Cref{prop:xi0-d4}, $\dps = 80$, spread $0$ across
$8$ representatives). Driver
\texttt{session\_q1\_newton.py} SHA-$256$
\texttt{1ad90f60$\dots$d5b10bc7}.

\item
\texttt{sessions/2026-05-03/Q20A-PHASE-C-RESUME/d2\_note\_v2/}.
The v2 LaTeX source frozen at concept-DOI deposit
(d2\_note\_v2.pdf SHA-$256$
\texttt{b9954d12$\dots$ece66}).

\item
\texttt{sessions/2026-05-03/D2-NOTE-V2-1-WASOW-FULL-CLOSURE/}.
This session: v2.1 build artefacts, Phase~B Newton-polygon
Lemma, Phase~C Birkhoff--Trjitzinsky~1933~\S\S4--6
verification, full-closure synthesis, claims.jsonl,
Zenodo runbook for the new-version deposit on concept
DOI \texttt{10.5281/zenodo.19996689}.
\end{itemize}

\bibliographystyle{plain}
\bibliography{annotated_bibliography}

\end{document}
